\appendix

\Needspace{16\baselineskip}
\section{The analytic facts used by the proof}
\label{app:analytic-framework}

The main argument uses the following analytic ingredients at several
derivative orders.  We collect them here so that every Sobolev, multiplier,
trace, and local-construction step invoked in the body is proved or attached
to a precise source.

\subsection{Fourier convention, Sobolev spaces, and frequency cutoffs}

For \(f\in\mathcal S(\R^3)\), set
\[
 \widehat f(\xi)
 :=(2\pi)^{-3/2}\int_{\R^3}e^{-ix\cdot\xi}f(x)\,dx,
 \qquad
 \langle\xi\rangle:=(1+|\xi|^2)^{1/2}.
\]
With this normalization, Plancherel's theorem is an isometry.
For every real \(s\), define
\[
 \norm{f}_{H^s}^2
 :=\int_{\R^3}\langle\xi\rangle^{2s}|\widehat f(\xi)|^2\,d\xi.
\]
The completion defines \(H^s(\R^3)\), and \(H^s_\sigma\) is its closed divergence-free subspace.
All multipliers below act componentwise on vector fields, and repeated spatial indices are summed.
For \(s_1\ge s_0\), the pointwise inequality between Fourier weights gives the continuous embedding
\[
 H^{s_1}\hookrightarrow H^{s_0}.
\]

The homogeneous operator \(\Lambda=(-\Delta)^{1/2}\) has symbol \(|\xi|\).
If \(0\le\theta\le k\), then
\[
 \norm{\Lambda^\theta f}_2^2
 =\int_{\R^3}|\xi|^{2\theta}|\widehat f(\xi)|^2\,d\xi
 \le\norm{f}_{H^k}^2.
\]
This comparison is what allows an inhomogeneous rectangle to control each homogeneous half-integer energy used in Part~II.
Every homogeneous quantity in the paper is evaluated on a function already known to lie in the indicated inhomogeneous Sobolev space, so no quotient by polynomials enters the argument.

The all-orders space is
\[
 H^\infty(\R^3):=\bigcap_{m\in\Nzero}H^m(\R^3),
 \qquad
 H^\infty_\sigma:=H^\infty\cap H^0_\sigma.
\]
Its topology is generated by the seminorms \(f\mapsto\norm{f}_{H^m}\).
Thus \(f_j\to f\) in \(H^\infty\) precisely when convergence holds in each fixed \(H^m\).
No sum over \(m\), and no constant uniform in \(m\), is part of this definition.

Choose a real, radial function \(\chi\in C_c^\infty(\R^3)\) with \(\chi=1\) on \(\{|\xi|\le1\}\) and \(\chi=0\) on \(\{|\xi|\ge2\}\), and define the low-frequency cutoff \(S_R\) by
\[
 \widehat{S_Rf}(\xi):=\chi(\xi/R)\widehat f(\xi).
\]
Dominated convergence gives \(S_Rf\to f\) strongly in every \(H^s\).
The same cutoffs act boundedly on every Bochner space \(L^q(J;H^s)\) with \(1\le q<\infty\), where the strong convergence follows first on simple functions and then by density.
These properties justify the frequency truncations used in the critical-height and endpoint-trace calculations.

\Needspace{16\baselineskip}
\subsection{Dyadic shells and the Sobolev weight they carry}

Part~II resolves the critical height into dyadic shells, so the relation between a shell amplitude and a Sobolev norm has to be exact up to fixed partition constants.
Let
\[
 \varphi(\xi):=\chi(\xi)-\chi(2\xi),
 \qquad
 \varphi_j(\xi):=\varphi(2^{-j}\xi),
 \qquad
 \widehat{P_jf}:=\varphi_j\widehat f.
\]
For \(\xi\ne0\), the telescoping sum satisfies
\[
 \sum_{j\in\mathbb Z}\varphi_j(\xi)=1.
\]
Each \(\varphi_j\) is supported where \(2^{j-1}\le|\xi|\le2^{j+1}\), and any fixed frequency belongs to only finitely many of these supports.

\begin{lemma}[Dyadic Sobolev equivalence]
For every \(s\in\R\), there are constants \(0<c_s\le C_s<\infty\), depending only on \(s\) and the chosen partition, such that
\begin{equation}
 c_s\norm{\Lambda^sf}_2^2
 \le\sum_{j\in\mathbb Z}2^{2js}\norm{P_jf}_2^2
 \le C_s\norm{\Lambda^sf}_2^2
 \label{eq:appendix-dyadic-equivalence}
\end{equation}
whenever either side is finite.
If \(K=2^j\) and \(A_K(f):=\frac12\norm{P_jf}_2^2\), then
\[
 \frac12\norm{\Lambda^sf}_2^2\simeq_s\sum_KK^{2s}A_K(f).
\]
\end{lemma}

\begin{proof}
Plancherel gives
\[
 \sum_j2^{2js}\norm{P_jf}_2^2
 =\int_{\R^3}
 \left(\sum_j2^{2js}|\varphi_j(\xi)|^2\right)
 |\widehat f(\xi)|^2\,d\xi.
\]
On the support of \(\varphi_j\), the quantities \(2^j\) and \(|\xi|\) differ by a factor between \(1/2\) and \(2\), so \(2^{2js}\simeq_s|\xi|^{2s}\).
Finite overlap gives the upper bound in \eqref{eq:appendix-dyadic-equivalence}.
For the lower bound, Cauchy--Schwarz and the uniform overlap number give
\[
 1=\left|\sum_j\varphi_j(\xi)\right|^2
 \le C\sum_j|\varphi_j(\xi)|^2
 \qquad(\xi\ne0).
\]
Multiplication by the comparable weights yields
\[
 |\xi|^{2s}\le C_s\sum_j2^{2js}|\varphi_j(\xi)|^2.
\]
Integration against \(|\widehat f|^2\) proves the lower bound.
The single frequency \(\xi=0\) has Lebesgue measure zero and contributes nothing to the \(L^2\) identity.
\end{proof}

For the critical height \(H=\frac12\norm{\Lambda^{1/2}u}_2^2\), the lemma gives
\[
 H(t)\simeq\sum_KK A_K(u(t)).
\]
At the next half-integer height, the same shell amplitude carries the weight \(K^3\), and each further full spatial derivative multiplies that weight by \(K^2\).
This is the precise frequency statement behind the vertical movement through the critical-height tower.

\Needspace{15\baselineskip}
\subsection{Spatial embeddings and interpolation}

The whole-space Sobolev embedding used throughout the proof follows directly from Fourier inversion.
If \(|\alpha|\le k\), then
\begin{align*}
 \norm{\partial^\alpha f}_\infty
 &\le (2\pi)^{-3/2}\int_{\R^3}|\xi|^{|\alpha|}|\widehat f(\xi)|\,d\xi
 \\
 &\le C\norm{f}_{H^s}
 \left(\int_{\R^3}|\xi|^{2|\alpha|}\langle\xi\rangle^{-2s}\,d\xi\right)^{1/2}.
\end{align*}
The last integral is finite under the strict inequality \(s>|\alpha|+3/2\).
Consequently,
\begin{equation}
 H^s(\R^3)\hookrightarrow C_b^k(\R^3)
 \qquad\text{whenever}\qquad s>k+\frac32.
 \label{eq:appendix-spatial-embedding}
\end{equation}
The finite-\(L^p\) version follows from the same dyadic decomposition.  Let
\(2\le p\le\infty\), \(|\alpha|\le k\), and
\(s>k+\frac32-\frac3p\).  Bernstein's inequality on the support of \(P_j\)
gives\bibref{BCD}
\[
 \norm{\partial^\alpha P_jf}_p
 \le C_{k,p}
 2^{j(k+3/2-3/p)}\norm{P_jf}_2
 \qquad(j\ge0).
\]
Put \(\delta=s-k-\frac32+\frac3p>0\).  Cauchy--Schwarz in \(j\) yields
\[
 \begin{aligned}
  \sum_{j\ge0}\norm{\partial^\alpha P_jf}_p
  &\le C_{k,p}
  \sum_{j\ge0}2^{-j\delta}
  \bigl(2^{js}\norm{P_jf}_2\bigr)\\
  &\le C_{s,k,p}\norm f_{H^s}.
 \end{aligned}
\]
The remaining low block \(S_{1/2}f\) obeys the same bound directly by
Bernstein's inequality.  Thus
\begin{equation}
 H^s(\R^3)\hookrightarrow W^{k,p}(\R^3)
 \qquad\text{whenever}\qquad
 2\le p\le\infty,
 \quad s>k+\frac32-\frac3p.
 \label{eq:appendix-Wkp-embedding}
\end{equation}
The strict condition is a single sufficient range for every \(p\); no endpoint
claim is needed in the body.
In particular,
\[
 H^3(\R^3)\hookrightarrow W^{1,\infty}(\R^3).
\]
This embedding turns the upper coordinate of the first concrete rectangle into an integrable Lipschitz coefficient, and it makes the local lifespan depend only on the base \(H^3\) norm.

Two homogeneous estimates drive the first differentiated energy calculation.
For \(f\in C_c^\infty(\R^3)\), the whole-space Sobolev inequality and Hölder interpolation give\bibref{Lieb}
\begin{equation}
 \norm{f}_6\le C\norm{\nabla f}_2,
 \qquad
 \norm{f}_r\le\norm{f}_2^\theta\norm{f}_6^{1-\theta},
 \qquad
 \frac1r=\frac\theta2+\frac{1-\theta}{6},
 \qquad
 2\le r\le6.
 \label{eq:appendix-homogeneous-Sobolev-interpolation}
\end{equation}
The interpolation estimate follows by writing \(|f|^r=(|f|^2)^{\theta r/2}(|f|^6)^{(1-\theta)r/6}\) and applying Hölder with the conjugate exponents forced by the displayed relation.
Density extends both estimates to the spaces in which they are used.
Applied componentwise to \(\nabla u\), the Sobolev inequality gives
\[
 \norm{\nabla u}_6\le C\norm{\nabla^2u}_2=C\norm{\Delta u}_2,
\]
where the equality follows from Plancherel after summing all second derivatives.

The critical response uses the homogeneous endpoint
\[
 \norm{f}_{L^3(\R^3)}\le C\norm{\Lambda^{1/2}f}_2.
\]
For Schwartz functions this is the Hardy--Littlewood--Sobolev inequality applied to \(f=\Lambda^{-1/2}(\Lambda^{1/2}f)\), and density gives the stated form\bibref{Lieb}.
Applying it to \(u\), \(\nabla u\), and \(\Lambda u\) gives
\[
 \norm{u}_3\le C\norm{\Lambda^{1/2}u}_2,
 \qquad
 \norm{\nabla u}_3+\norm{\Lambda u}_3\le C\norm{\Lambda^{3/2}u}_2.
\]
These are the three factors used in the critical-response estimate \eqref{eq:critical-response-estimate}.

\Needspace{15\baselineskip}
\subsection{Products, transport, and commutators}

The multiplication estimates can be read directly from the Fourier convolution formula.
For \(s\ge0\), the elementary weight inequality
\[
 \langle\xi\rangle^s\le C_s\bigl(\langle\eta\rangle^s+\langle\xi-\eta\rangle^s\bigr)
\]
and Young's convolution inequality give
\[
 \norm{fg}_{H^s}
 \le C_s\left(
 \norm{\langle\cdot\rangle^s\widehat f}_2\norm{\widehat g}_1
 +\norm{\widehat f}_1\norm{\langle\cdot\rangle^s\widehat g}_2
 \right).
\]
When \(s>3/2\), Cauchy--Schwarz gives \(\norm{\widehat f}_1\le C_s\norm f_{H^s}\), because \(\langle\xi\rangle^{-s}\in L^2(\R^3)\).
Hence
\begin{equation}
 \norm{fg}_{H^s}\le C_s\norm f_{H^s}\norm g_{H^s},
 \qquad
 s>\frac32.
 \label{eq:appendix-Hs-algebra}
\end{equation}
This is the inhomogeneous Sobolev-algebra estimate; its fractional and unequal-order variants follow from the same Littlewood--Paley calculus\bibref{BCD}.

The lowest rectangle requires a separate estimate because \(H^1(\R^3)\) is below the algebra threshold.
Interpolation between \(L^2\) and \(L^6\) gives \(H^1\hookrightarrow L^4\), so Hölder gives
\begin{equation}
 \norm{fg}_2\le\norm f_4\norm g_4\le C\norm f_{H^1}\norm g_{H^1}.
 \label{eq:appendix-low-rectangle-product}
\end{equation}
For every integer \(M\ge1\), the algebra estimate at order \(M+1>3/2\) gives
\[
 \norm{fg}_{H^M}\le C_M\norm f_{H^{M+1}}\norm g_{H^{M+1}}.
\]
Together with \eqref{eq:appendix-low-rectangle-product}, this proves for all \(M\in\Nzero\) the rectangle product law
\begin{equation*}
 \norm{fg}_{H^M}\le C_M\norm f_{H^{M+1}}\norm g_{H^{M+1}}.
\end{equation*}
The two cases have different proofs, and keeping them separate prevents an algebra property from being attributed to \(H^1\).

Taking one more derivative before using the algebra gives the complete-row transport estimate
\[
 \norm{\PP\nabla\cdot(f\otimes g)}_{H^m}
 \le C_m\norm f_{H^{m+2}}\norm g_{H^{m+2}},
 \qquad
 m\in\Nzero.
\]
For the base local construction, the sharper placement is
\begin{equation}
 \norm{\PP\nabla\cdot(f\otimes g)}_{H^2}
 \le C\norm f_{H^3}\norm g_{H^3}.
 \label{eq:appendix-local-product}
\end{equation}
Indeed, \eqref{eq:appendix-Hs-algebra} at \(s=3\) controls \(f\otimes g\) in \(H^3\), after which divergence lowers the order by one and the Leray projection preserves it.

For higher-order energy estimates, the undifferentiated transport term must be kept apart from the commutator containing the differentiated coefficients.
Write \(\mathcal J^m=(1-\Delta)^{m/2}\).
The Kato--Ponce commutator theorem gives, for every integer \(m\ge3\),
\begin{equation}
 \norm{[\mathcal J^m,f\cdot\nabla]g}_2
 \le C_m\left(
 \norm{\nabla f}_\infty\norm g_{H^m}
 +\norm f_{H^m}\norm{\nabla g}_\infty
 \right).
 \label{eq:appendix-commutator}
\end{equation}
This is the exact specialization of the standard commutator estimate used here\bibref{KatoPonce}.
When \(f=g=w\), it becomes
\[
 \norm{[\mathcal J^m,w\cdot\nabla]w}_2
 \le C_m\norm{\nabla w}_\infty\norm w_{H^m}.
\]
The integer multi-index form used in Section~\ref{sec:high-order-continuation} is the same estimate before the derivatives are collected:
\[
 \sum_{|\alpha|\le m}
 \left|\bigl([\partial^\alpha,w\cdot\nabla]w,\partial^\alpha w\bigr)_2\right|
 \le C_m\norm{\nabla w}_\infty\norm w_{H^m}^2.
\]
The embedding \eqref{eq:appendix-spatial-embedding} supplies the \(L^\infty\) factors in every application made in the paper.

\Needspace{16\baselineskip}
\subsection{Riesz transforms, the Leray projection, and normalized pressure}

For \(\xi\ne0\), the Riesz transform \(R_j\) has symbol \(i\xi_j/|\xi|\), and the Leray projection has matrix symbol
\[
 \widehat{\PP}(\xi)
 =I-\frac{\xi\otimes\xi}{|\xi|^2}.
\]
Their values at \(\xi=0\) may be fixed arbitrarily.
Since these symbols are bounded, Plancherel gives, for every \(s\in\R\),
\begin{equation}
 \norm{R_jf}_{H^s}\le\norm f_{H^s},
 \qquad
 \norm{\PP f}_{H^s}\le C\norm f_{H^s}.
 \label{eq:appendix-riesz-multiplier-bounds}
\end{equation}
The same bounds hold for the vector and tensor fields used below.
These are the Sobolev instances of the order-zero multiplier theorem\bibref{Stein}.

The pressure formula, including its sign, follows from the unprojected momentum equation.
If \(\nabla\cdot u=0\), then
\[
 (u\cdot\nabla)u_i=\partial_j(u_ju_i).
\]
Taking the divergence cancels the time and viscous terms and gives
\begin{equation}
 -\Delta p=\partial_i\partial_j(u_i u_j).
 \label{eq:appendix-pressure-poisson}
\end{equation}
The symbol of \(R_iR_j\) is \(-\xi_i\xi_j/|\xi|^2\), so
\begin{equation}
 p=R_iR_j(u_i u_j)
 \label{eq:appendix-pressure-riesz}
\end{equation}
solves \eqref{eq:appendix-pressure-poisson}.
If two such pressures belong to \(L^2(\R^3)\), their difference \(h\) is a harmonic \(L^2\) function.
The distributional identity \(|\xi|^2\widehat h=0\) says that the \(L^2\) function \(\widehat h\) is supported at the single point \(\xi=0\), and hence \(h=0\).
Requiring the pressure to belong to any nonnegative inhomogeneous Sobolev space consequently fixes both the spatial constant and any time-dependent constant.

Let \(F=\nabla\cdot(u\otimes u)\).
Because \(-\Delta p=\nabla\cdot F\), the Fourier symbols of the two sides give
\begin{equation}
 \PP F=F+\nabla p.
 \label{eq:appendix-projected-pressure}
\end{equation}
Thus the projected equation and the normalized pressure-complete equation are equivalent in every Sobolev class used in the paper.
Temporal and spatial derivatives commute with the Riesz transforms and with \(\PP\), so the same normalization passes through every row of the mixed jet.

\Needspace{15\baselineskip}
\subsection{The two Leibniz rules behind the mixed jet}

The temporal recurrence and the mixed spatial recurrence are generated by ordinary differentiation of one product.
Writing both combinatorial factors once prevents either factor from being hidden when a pressure or transport row is used later.

\begin{lemma}[Mixed Leibniz identity]
\label{lem:appendix-mixed-leibniz}
Let \(n\in\Nzero\), let \(\alpha\) be a spatial multi-index, and let \(f,g\) be smooth enough for the derivatives below to exist.
Then
\begin{equation}
 \partial_t^n\partial_x^\alpha(fg)
 =
 \sum_{a=0}^{n}\sum_{\beta\le\alpha}
 \binom na\binom\alpha\beta
 \bigl(\partial_t^a\partial_x^\beta f\bigr)
 \bigl(\partial_t^{n-a}\partial_x^{\alpha-\beta}g\bigr).
 \label{eq:appendix-mixed-leibniz}
\end{equation}
The same identity holds for the Sobolev rows used in this paper.
\end{lemma}

\begin{proof}
The one-variable Leibniz rule gives
\[
 \partial_t^n(fg)=\sum_{a=0}^n\binom na(\partial_t^af)(\partial_t^{n-a}g).
\]
For each summand, the multi-index spatial rule gives
\[
 \partial_x^\alpha(FG)
 =\sum_{\beta\le\alpha}\binom\alpha\beta
 (\partial_x^\beta F)(\partial_x^{\alpha-\beta}G).
\]
Spatial and temporal derivatives commute, so applying the second identity to every term in the first produces \eqref{eq:appendix-mixed-leibniz}.
For the Sobolev rows used here, spatial smoothing gives the classical identity, and the product estimates above pass every term to the limit in the fixed target topology.
\end{proof}

Applied to the product in \eqref{eq:appendix-pressure-riesz} and followed by \(R_iR_j\), the lemma gives
\[
 \partial_t^n\partial_x^\alpha p
 =
 R_iR_j
 \sum_{a=0}^{n}\sum_{\beta\le\alpha}
 \binom na\binom\alpha\beta
 \bigl(\partial_t^a\partial_x^\beta u_i\bigr)
 \bigl(\partial_t^{n-a}\partial_x^{\alpha-\beta}u_j\bigr),
\]
which is the formula used for the pressure coordinates in \eqref{eq:pressure-time-jet} and \eqref{eq:mixed-pressure}.
Applied to \(u\cdot\nabla u\), it gives the corresponding transport rows.
Every term at a fixed \((n,\alpha)\) uses only temporal orders at most \(n\) and spatial orders at most \(|\alpha|+1\), which is the finite dependency count required by each rectangle.

\Needspace{17\baselineskip}
\section{Time regularity beyond the middle trace}
\label{app:time-traces}

Lemma~\ref{lem:part-ii-adjacent-trace} already proves that the two outer rows
\(L^2H^{s+1}\) and \(L^2H^{s-1}\) meet in one strong \(H^s\) trace.  This
appendix records the shifted pairing behind that argument, the
Banach-valued fundamental theorem of calculus used to identify time
derivatives, and the higher time-Sobolev embedding used for the complete jet.

\subsection{Shifted duality}

The Hilbert triple is
\[
 H^{s+1}\hookrightarrow H^s\hookrightarrow H^{s-1}.
\]
For \(f\in H^{s-1}\) and \(g\in H^{s+1}\), define the outer-space pairing with \(H^s\) as pivot by
\[
 \pair{f}{g}_{H^{s-1},H^{s+1};H^s}
 :=\int_{\R^3}\langle\xi\rangle^{2s}
 \widehat f(\xi)\cdot\overline{\widehat g(\xi)}\,d\xi.
\]
The exact factorization
\[
 \langle\xi\rangle^{2s}
 =\langle\xi\rangle^{s-1}\langle\xi\rangle^{s+1}
\]
and Cauchy--Schwarz give
\begin{equation}
 \left|\pair{f}{g}_{H^{s-1},H^{s+1};H^s}\right|
 \le\norm f_{H^{s-1}}\norm g_{H^{s+1}}.
 \label{eq:appendix-shifted-duality-bound}
\end{equation}
This pairing is the exact one used both in the terminal trace and in the linear \(H^3\) estimate of Appendix~\ref{app:local-theory}.

\subsection{The Banach-valued fundamental theorem of calculus}

The time-jet construction first identifies a derivative in distributions and then needs to know that the corresponding row is a genuine time-differentiable curve.
The following elementary result is the bridge.

\begin{lemma}[Banach-valued fundamental theorem of calculus]
\label{lem:banach-fundamental-calculus}
Let \(X\) be a Banach space, let \(J=(a,b)\) be a bounded interval, and suppose \(v,g\in L^1(J;X)\) satisfy \(\partial_tv=g\) in distributions.
Then \(v\) has a unique absolutely continuous representative on \([a,b]\), still denoted by \(v\), and
\begin{equation}
 v(t)-v(r)=\int_r^tg(s)\,ds
 \qquad(a\le r\le t\le b).
 \label{eq:appendix-banach-ftc}
\end{equation}
In particular, if \(g\in C([a,b];X)\), then \(v\in C^1([a,b];X)\) and its classical derivative is \(g\).
\end{lemma}

\begin{proof}
Fix \(t_0\in J\) and define the Bochner integral
\[
 G(t):=\int_{t_0}^tg(s)\,ds.
\]
The map \(G\) is absolutely continuous, belongs to \(L^1(J;X)\), and has distributional derivative \(g\).
Thus \(h:=v-G\) belongs to \(L^1(J;X)\) and has distributional derivative zero.

\Needspace{9\baselineskip}
Let \(\rho_\varepsilon\) be a scalar time mollifier and set \(h_\varepsilon=\rho_\varepsilon*h\) on the interior interval \(J_\varepsilon=(a+\varepsilon,b-\varepsilon)\).
Convolution commutes with the distributional derivative, so \((h_\varepsilon)'=0\) and \(h_\varepsilon\) is an \(X\)-valued constant on \(J_\varepsilon\).
On any compact interval \(K\Subset J\), the functions \(h_\varepsilon\) converge to \(h\) in \(L^1(K;X)\).
Since each \(h_\varepsilon\) is constant on \(K\), the identity
\[
 |K|\,\norm{h_\varepsilon-h_\delta}_X
 =\norm{h_\varepsilon-h_\delta}_{L^1(K;X)}
\]
shows that these constants form a Cauchy family.
Their limit is a vector \(c_K\in X\), and \(h=c_K\) almost everywhere on \(K\).
Overlapping compact intervals force the constants to agree, so one vector \(c\in X\) satisfies \(h=c\) almost everywhere on \(J\).
Consequently, \(v=c+G\) almost everywhere, and this formula supplies the absolutely continuous representative and identity \eqref{eq:appendix-banach-ftc}.
Two continuous representatives agreeing almost everywhere agree everywhere, which proves uniqueness.
If \(g\) is continuous, the ordinary differentiation theorem for the Bochner integral gives \(G'=g\) at every time.
\end{proof}

\subsection{Higher time regularity}

After the compatible rectangles have been intersected, the paper also uses the ordinary embedding of a finite time-Sobolev row into classical time derivatives.

\begin{lemma}[Time-Sobolev embedding with values in a Hilbert space]
Let \(X\) be a Hilbert space, let \(r\ge1\) be an integer, and define
\[
 H^r(J;X)
 :=
 \left\{v\in L^2(J;X):\partial_t^jv\in L^2(J;X),\ 1\le j\le r\right\}.
\]
If \(a\in\Nzero\) and \(r>a+\frac12\), then
\begin{equation}
 H^r(J;X)\hookrightarrow C^a(\overline J;X).
 \label{eq:appendix-time-embedding}
\end{equation}
\end{lemma}

\begin{proof}
Since \(r\) and \(a\) are integers, the strict inequality says \(a\le r-1\).
For each \(0\le j\le a\), the functions \(\partial_t^jv\) and \(\partial_t^{j+1}v\) belong to \(L^2(J;X)\), and hence to \(L^1(J;X)\) because \(J\) is bounded.
Lemma~\ref{lem:banach-fundamental-calculus} gives a continuous representative of \(\partial_t^jv\), and it identifies the derivative of the \(j\)-th representative with the \((j+1)\)-st.
Induction from \(j=0\) through \(j=a\) yields \(v\in C^a(\overline J;X)\).
The embedding is continuous by the closed graph theorem, or directly from \eqref{eq:appendix-banach-ftc} after choosing a time at which the pointwise norm is bounded by the \(L^2\) average.
\end{proof}

Taking \(X=H^m(\R^3)\) in \eqref{eq:appendix-time-embedding} and then applying \eqref{eq:appendix-spatial-embedding} converts the all-orders Bochner control into the mixed classical regularity used in the endpoint and local-jet arguments.

\Needspace{18\baselineskip}
\section{Local construction, pressure recovery, and restart}
\label{app:local-theory}

The endpoint argument has to restart the same pressure-normalized equation from the terminal velocity.
This section uses the time tools collected in Appendix~\ref{app:time-traces} to construct that local solution in the base space \(H^3_\sigma\), prove that its lifespan depends only on the base norm and the viscosity, and retain every higher order already present in the datum.
The result is the local theorem stated as Proposition~\ref{prop:maximal-history}.

\subsection{\texorpdfstring{The linear \(H^3\) estimate}{The linear H3 estimate}}

Let \(J_\delta=[t_0,t_0+\delta]\) with \(0<\delta\le1\), and consider
\[
 z_t-\nu\Delta z=F,
 \qquad
 z(t_0)=a,
\]
where \(a\in H^3\) and \(F\in L^2(J_\delta;H^2)\).
For a smooth solution, take the \(H^3\) energy pairing and place the force in the shifted \(H^2,H^4\) duality from \eqref{eq:appendix-shifted-duality-bound}.
This gives
\begin{equation}
 \frac12\frac d{dt}\norm z_{H^3}^2
 +\nu\norm{\nabla z}_{H^3}^2
 =\operatorname{Re}\pair{F}{z}_{H^2,H^4;H^3}.
 \label{eq:local-linear-energy}
\end{equation}
The inhomogeneous Fourier weights satisfy the exact identity
\begin{equation}
 \norm z_{H^4}^2
 =\norm z_{H^3}^2+\norm{\nabla z}_{H^3}^2.
 \label{eq:local-weight-splitting}
\end{equation}
Using shifted duality, \eqref{eq:local-weight-splitting}, and Young's inequality in \eqref{eq:local-linear-energy} gives
\[
 \frac d{dt}\norm z_{H^3}^2
 +\nu\norm{\nabla z}_{H^3}^2
 \le C_\nu\left(\norm z_{H^3}^2+\norm F_{H^2}^2\right).
\]
Gronwall's inequality, time integration, and \eqref{eq:local-weight-splitting} now yield
\[
 \norm z_{C(J_\delta;H^3)}
 +\norm z_{L^2(J_\delta;H^4)}
 \le C_\nu\left(\norm a_{H^3}+\norm F_{L^2(J_\delta;H^2)}\right).
\]
The equation gives \(z_t=\nu\Delta z+F\in L^2(J_\delta;H^2)\), and hence
\begin{equation}
 \norm z_{C(J_\delta;H^3)}
 +\norm z_{L^2(J_\delta;H^4)}
 +\norm{z_t}_{L^2(J_\delta;H^2)}
 \le C_\nu\left(\norm a_{H^3}+\norm F_{L^2(J_\delta;H^2)}\right).
 \label{eq:local-maximal-regularity}
\end{equation}
The constant is uniform for \(0<\delta\le1\).  Its viscosity dependence can
be kept explicit.  Define
\[
 L_1(\nu):=e^{\nu/2}\max\{1,\nu^{-1/2}\},
\]
\[
 L_\nabla(\nu)
 :=\max\left\{
 \left(\frac{1+\nu e^\nu}{\nu}\right)^{1/2},
 \left(\frac{e^\nu+\nu^{-1}}{\nu}\right)^{1/2}
 \right\},
\]
and
\begin{equation}
 C_{\rm lin}(\nu)
 :=1+L_1(\nu)
 +(1+\nu)\bigl(L_1(\nu)^2+L_\nabla(\nu)^2\bigr)^{1/2}.
 \label{eq:explicit-linear-local-constant}
\end{equation}
Indeed, Young's inequality in \eqref{eq:local-linear-energy} gives
\[
 \frac d{dt}\norm z_{H^3}^2
 +\nu\norm{\nabla z}_{H^3}^2
 \le\nu\norm z_{H^3}^2+\nu^{-1}\norm F_{H^2}^2.
\]
Gronwall on an interval of length at most one gives the \(L_1(\nu)\)
coefficient for the uniform \(H^3\) norm.  Integration gives the
\(L_\nabla(\nu)\) coefficient for the gradient term.  The exact splitting
\eqref{eq:local-weight-splitting} and
\(z_t=\nu\Delta z+F\) then show that \(C_{\rm lin}(\nu)\) is one admissible
constant on the right of \eqref{eq:local-maximal-regularity}.

For completeness, this estimate also constructs the nonsmooth solution.
Apply the cutoff \(S_R\) to the datum and force, solve the resulting frequency-truncated equation, and use \eqref{eq:local-maximal-regularity} on differences.
The truncated solutions are Cauchy in all three spaces on the left of \eqref{eq:local-maximal-regularity}.
Their limit solves the equation in distributions, and Lemma~\ref{lem:part-ii-adjacent-trace} with \(m=3\) supplies the strong \(H^3\) trace and the value \(z(t_0)=a\).
Uniqueness follows by applying the same estimate to the difference of two solutions.

Taking the spatial Fourier transform and solving the scalar ordinary differential equation at each frequency gives the equivalent Duhamel formula
\[
 z(t)=e^{\nu(t-t_0)\Delta}a
 +\int_{t_0}^{t}e^{\nu(t-s)\Delta}F(s)\,ds.
\]
The plus sign records the sign of the force \(F\); the Navier--Stokes force inserted below is negative.

\Needspace{18\baselineskip}
\subsection{The nonlinear fixed point and the base lifespan}

Applying the Leray projection to the velocity equation gives the mild map
\begin{equation}
 \Phi(w)(t)
 :=e^{\nu(t-t_0)\Delta}a
 -\int_{t_0}^{t}e^{\nu(t-s)\Delta}
 \PP\nabla\cdot(w\otimes w)(s)\,ds.
 \label{eq:local-mild}
\end{equation}
Set
\[
 X_\delta
 :=C(J_\delta;H^3_\sigma)\cap L^2(J_\delta;H^4),
 \qquad
 \norm w_{X_\delta}
 :=\norm w_{C H^3}+\norm w_{L^2H^4}.
\]
All abbreviated time norms in this subsection are taken over \(J_\delta\).
Let \(C_{\rm quad}\) be a valid dimension-dependent constant in the local
product estimate \eqref{eq:appendix-local-product}.  It gives
\begin{equation}
 \norm{\PP\nabla\cdot(w\otimes w)}_{L^2H^2}
 \le C_{\rm quad}\delta^{1/2}\norm w_{C H^3}^2,
 \label{eq:local-index-count}
\end{equation}
and the decomposition
\[
 w\otimes w-v\otimes v
 =(w-v)\otimes w+v\otimes(w-v)
\]
gives
\begin{equation}
 \norm{\PP\nabla\cdot(w\otimes w-v\otimes v)}_{L^2H^2}
 \le C_{\rm quad}\delta^{1/2}
 \left(\norm w_{C H^3}+\norm v_{C H^3}\right)
 \norm{w-v}_{C H^3}.
 \label{eq:local-difference-count}
\end{equation}

Estimate \eqref{eq:local-maximal-regularity} applied to \eqref{eq:local-mild} yields
\[
 \norm{\Phi(w)}_{X_\delta}
 \le C_{\rm lin}(\nu)\norm a_{H^3}
 +C_{\rm lin}(\nu)C_{\rm quad}\delta^{1/2}
  \norm w_{X_\delta}^2
\]
and
\[
 \norm{\Phi(w)-\Phi(v)}_{X_\delta}
 \le C_{\rm lin}(\nu)C_{\rm quad}\delta^{1/2}
 \left(\norm w_{X_\delta}+\norm v_{X_\delta}\right)
 \norm{w-v}_{X_\delta}.
\]
If \(a=0\), the zero solution is the unique fixed point.
For \(a\ne0\), take
\(R=2C_{\rm lin}(\nu)\norm a_{H^3}\) and choose \(\delta\le1\) so that
\[
 4C_{\rm lin}(\nu)C_{\rm quad}\delta^{1/2}R\le1.
\]
Then \(\Phi\) maps the closed \(X_\delta\)-ball of radius \(R\) into itself and has Lipschitz constant at most \(1/2\) there.
Banach's fixed-point theorem produces a solution \(w\in X_\delta\).
The choice above permits the quantitative lower bound
\begin{equation}
 \delta(\nu,\norm a_{H^3})
 \ge\min\left\{1,c_\nu(1+\norm a_{H^3})^{-2}\right\}.
 \label{eq:local-lifespan}
\end{equation}
Here one may take
\begin{equation}
 c_\nu
 :=\left(
 8C_{\rm quad}C_{\rm lin}(\nu)^2
 \right)^{-2}.
 \label{eq:explicit-local-lifespan-constant}
\end{equation}
Only the \(H^3\) size of the datum and the fixed viscosity enter this lifespan.

The fixed point is unique throughout \(X_\delta\).
If \(w,v\in X_\delta\) have the same datum, divide \(J_\delta\) into finitely many intervals whose lengths make the coefficient obtained from \eqref{eq:local-difference-count} and \eqref{eq:local-maximal-regularity} smaller than one.
The difference vanishes on the first interval.
Its common endpoint becomes zero initial data on the next interval, and finite induction gives \(w=v\) throughout \(J_\delta\).
The heat semigroup and \(\PP\) preserve divergence-free fields, so \(w(t)\in H^3_\sigma\).
The projected equation and \eqref{eq:local-index-count} also give
\[
 w_t\in L^2(J_\delta;H^2),
\]
which completes the base regularity in \eqref{eq:local-base-class}.

\subsection{Recovering the normalized pressure}

The fixed point satisfies
\[
 w_t+\PP\nabla\cdot(w\otimes w)=\nu\Delta w.
\]
Define
\[
 q:=R_iR_j(w_iw_j).
\]
Because \(w\in C(J_\delta;H^3)\) and \(H^3\) is an algebra, the Riesz bounds give \(q\in C(J_\delta;H^3)\).
With \(F=\nabla\cdot(w\otimes w)\), identity \eqref{eq:appendix-projected-pressure} gives
\[
 \PP F=F+\nabla q.
\]
Thus the projected fixed point solves
\[
 w_t+(w\cdot\nabla)w+\nabla q=\nu\Delta w,
 \qquad
 \nabla\cdot w=0.
\]
The condition \(q(t)\in H^3(\R^3)\) fixes the whole-space normalization at each time.
Any second pressure in the same class has harmonic \(L^2\) difference and hence coincides with \(q\).
Uniqueness of the projected fixed point consequently gives uniqueness of the pressure-complete pair.

\Needspace{18\baselineskip}
\subsection{Persistence on the base lifespan}

Fix an integer \(m\ge3\) and first take smooth data.
Apply \(\mathcal J^m=(1-\Delta)^{m/2}\) to the projected equation and pair with \(\mathcal J^mw\).
The operators \(\mathcal J^m\), \(\PP\), and \(\Delta\) commute.
Since \(\mathcal J^mw\) is divergence-free,
\[
 \bigl((w\cdot\nabla)\mathcal J^mw,\mathcal J^mw\bigr)_2=0.
\]
After this cancelling transport term is separated, the energy identity is
\[
 \frac12\frac d{dt}\norm w_{H^m}^2
 +\nu\norm{\nabla w}_{H^m}^2
 =-\operatorname{Re}
 \bigl([\mathcal J^m,w\cdot\nabla]w,\mathcal J^mw\bigr)_2.
\]
The commutator estimate \eqref{eq:appendix-commutator} gives
\begin{equation}
 \frac12\frac d{dt}\norm w_{H^m}^2
 +\nu\norm{\nabla w}_{H^m}^2
 \le C_m\norm{\nabla w}_\infty\norm w_{H^m}^2.
 \label{eq:local-persistence}
\end{equation}
This is the local form of the higher-order continuation estimate \eqref{eq:high-order-continuation}.
The base solution belongs to \(C(J_\delta;H^3)\), and \eqref{eq:appendix-spatial-embedding} puts \(\nabla w\) in \(C(J_\delta;L^\infty)\).
Thus the coefficient in \eqref{eq:local-persistence} is integrable on the whole base interval.
Gronwall gives
\begin{equation}
 \norm{w(t)}_{H^m}^2
 \le\norm a_{H^m}^2
 \exp\left(
 2C_m\int_{t_0}^{t}\norm{\nabla w(s)}_\infty\,ds
 \right).
 \label{eq:local-persistence-bound}
\end{equation}
The dissipation term in \eqref{eq:local-persistence}, together with the weight splitting at order \(m\), gives
\[
 w\in L^2(J_\delta;H^{m+1}).
\]
The product estimate then puts \(\PP\nabla\cdot(w\otimes w)\) in \(L^2(J_\delta;H^{m-1})\), so the equation gives
\[
 w_t\in L^2(J_\delta;H^{m-1}).
\]
Lemma~\ref{lem:part-ii-adjacent-trace} supplies \(w\in C(J_\delta;H^m)\).

For a general datum \(a\in H^m_\sigma\), let \(a^{(k)}=S_ka\).
The data converge to \(a\) in \(H^m\), and their \(H^3\) norms are uniformly bounded.
The lifespan \eqref{eq:local-lifespan} hence places every smooth solution \(w^{(k)}\) on one common interval determined by the base norm.
The contraction estimate makes \(w^{(k)}\) converge to the \(H^3\) fixed point in \(C H^3\cap L^2H^4\).
Estimate \eqref{eq:local-persistence-bound} gives a uniform \(L^\infty H^m\) bound, while \eqref{eq:local-persistence} gives a uniform \(L^2H^{m+1}\) bound and the equation gives a uniform \(L^2H^{m-1}\) bound for the time derivative.
Weak compactness and the already known \(H^3\) limit identify all three weak limits with the same base solution \(w\).
Lemma~\ref{lem:part-ii-adjacent-trace} gives its \(H^m\)-continuous representative.
The energy inequality passed to the limit gives
\[
 \limsup_{t\downarrow t_0}\norm{w(t)}_{H^m}\le\norm a_{H^m},
\]
while weak convergence to the prescribed distributional trace gives the reverse lower bound.
Uniform convexity of \(H^m\) then yields \(w(t)\to a\) strongly in \(H^m\) as \(t\downarrow t_0\).
This completes persistence for nonsmooth \(H^m\) data without changing the base lifespan.

If \(a\in H^\infty_\sigma\), repeat the argument at each fixed \(m\).
The solutions obtained at different orders coincide by the \(H^3\) uniqueness already proved.
One solution consequently satisfies
\begin{equation}
 w\in C(J_\delta;H^m)\cap L^2(J_\delta;H^{m+1}),
 \qquad
 w_t\in L^2(J_\delta;H^{m-1})
 \qquad(m\in\N,\ m\ge3).
 \label{eq:local-all-spatial-orders}
\end{equation}
The constants may depend on \(m\), while the interval remains the one selected by the \(H^3\) datum.

\Needspace{18\baselineskip}
\subsection{The pressure-complete time jet}

All spatial orders in \eqref{eq:local-all-spatial-orders} allow the equation to generate temporal rows without leaving the construction interval.
Set \(V_0=w\).
Given rows \(V_0,\ldots,V_n\) that are continuous in every finite Sobolev space, define
\begin{equation}
 Q_n
 :=R_iR_j\sum_{a=0}^{n}\binom naV_{a,i}V_{n-a,j}
 \in C(J_\delta;H^m)
 \qquad(m\in\Nzero)
 \label{eq:local-pressure-time-row}
\end{equation}
and
\begin{equation}
 V_{n+1}
 :=\nu\Delta V_n
 -\sum_{a=0}^{n}\binom na(V_a\cdot\nabla)V_{n-a}
 -\nabla Q_n
 \in C(J_\delta;H^m)
 \qquad(m\in\Nzero).
 \label{eq:local-velocity-time-row}
\end{equation}
The membership statements follow from the product estimates above after choosing two extra spatial orders for diffusion and one extra order for transport and pressure.
They are asserted one target \(m\) at a time, so no constant uniform across the tower is required.

The recursive symbols are now identified with the actual derivatives.
At \(n=0\), the pressure-complete equation gives \(\partial_tV_0=V_1\) in distributions.
Since \(V_1\) is continuous in every target \(H^m\), Lemma~\ref{lem:banach-fundamental-calculus} gives \(V_0\in C^1(J_\delta;H^m)\) and identifies its classical derivative with \(V_1\).
Suppose inductively that \(V_j=\partial_t^jw\) for \(0\le j\le n\).
Differentiate the pressure formula and momentum equation \(n\) times in distributions.
The temporal instance of Lemma~\ref{lem:appendix-mixed-leibniz} produces the two binomial sums in \eqref{eq:local-pressure-time-row} and \eqref{eq:local-velocity-time-row}, hence \(\partial_tV_n=V_{n+1}\).
Another application of the Banach-valued fundamental theorem promotes that distributional derivative to a classical derivative in every fixed \(H^m\).
Induction gives
\[
 \partial_t^nw,\ \partial_t^nq
 \in C(J_\delta;H^m)
 \qquad(n,m\in\Nzero).
\]
This argument constructs derivatives at each finite order; it makes no assertion that the resulting formal Taylor series converges in time.

Evaluation at \(t=t_0\) shows that the complete right-hand jet is determined by the single datum \(a\):
\[
 V_0(t_0)=a,
\]
\[
 Q_n(t_0)
 =R_iR_j\sum_{b=0}^{n}\binom nb
 V_{b,i}(t_0)V_{n-b,j}(t_0),
\]
and
\[
 V_{n+1}(t_0)
 =\nu\Delta V_n(t_0)
 -\sum_{b=0}^{n}\binom nb
 (V_b(t_0)\cdot\nabla)V_{n-b}(t_0)
 -\nabla Q_n(t_0).
\]
These are the normalized right-hand recurrences compared with the terminal left jet in \eqref{eq:endpoint-pressure} and \eqref{eq:endpoint-velocity}.
The local construction starts with the same pressure gauge and the same equation-generated derivatives that arrive from the old interval.

\Needspace{17\baselineskip}
\subsection{\texorpdfstring{Maximal histories and the \(H^3\) alternative}{Maximal histories and the H3 alternative}}

Starting from \(t_0\), join successive local intervals whenever the current solution has an \(H^3\) endpoint value.
Uniqueness identifies the solutions on every overlap, so their union is one pressure-normalized history on a maximal interval \([t_0,T_{\max})\).

Assume \(T_{\max}<\infty\) and
\[
 \sup_{t_0\le t<T_{\max}}\norm{w(t)}_{H^3}\le M.
\]
The lifespan bound \eqref{eq:local-lifespan} supplies the same positive number
\[
 \delta_M:=\min\{1,c_\nu(1+M)^{-2}\}
\]
whenever the local theorem is restarted from a time \(t<T_{\max}\).
Choose \(t>T_{\max}-\delta_M/2\) and start the construction from the actual datum \(w(t)\).
The new local solution exists past \(T_{\max}\), and uniqueness identifies it with the maximal history on their overlap.
This contradicts maximality.
Hence
\[
 T_{\max}<\infty
 \quad\Longrightarrow\quad
 \limsup_{t\uparrow T_{\max}}\norm{w(t)}_{H^3}=\infty.
\]

If the initial datum belongs to \(H^m\) for some \(m\ge3\), the persistence argument keeps that order on every local piece of the maximal history.
A bounded \(H^m\) norm bounds the \(H^3\) norm, so the same late-slice restart proves the corresponding \(H^m\) alternative under its datum hypothesis.
The construction, pressure recovery, persistence, and continuation argument together prove Proposition~\ref{prop:maximal-history}.

At the endpoint in Part~IV, the compatible intersection supplies \(u_*\in H^\infty_\sigma\), and in particular \(u_*\in H^3_\sigma\).
The construction above produces a normalized solution on a right-hand interval whose length depends only on \(\nu\) and \(\norm{u_*}_{H^3}\).
The pressure-complete jet calculation shows that its right derivatives agree with the already reconstructed left jet.
The matching jets make the piecewise pair classical across \(T_*\), while local uniqueness preserves the original pair on its old interval.
